%Môn: Toán
\begin{ex}
    Cho cấp số nhân $(u_n)$ có $u_1 = 3$ và $u_2 = 6$. Tìm công bội $q$.
    \shortans{2}
    \loigiai{Công bội $q = u_2 / u_1 = 6 / 3 = 2$.}
\end{ex}

%Môn: Vật Lí
\begin{ex}
    Một vật dao động điều hòa với tần số góc $\omega = 10\pi$ rad/s. Chu kì dao động $T$ là bao nhiêu giây?
    \shortans{0,2}
    \loigiai{$T = 2\pi / \omega = 2\pi / 10\pi = 0,2$ s.}
\end{ex}

%Môn: Hóa Học
\begin{ex}
    Tính giá trị pH của dung dịch có nồng độ ion $H^+$ bằng $0,01$ M.
    \shortans{2}
    \loigiai{$pH = -\log(0,01) = 2$.}
\end{ex}

%Môn: Địa Lí
\begin{ex}
    Trong các quốc gia sau, quốc gia nào có diện tích lớn nhất thế giới?
    \choice
    {Hoa Kỳ}
    {Trung Quốc}
    {\True Liên bang Nga}
    {Canada}
    \loigiai{Liên bang Nga có diện tích khoảng 17 triệu $km^2$, đứng đầu thế giới.}
\end{ex}

%Môn: Lịch Sử
\begin{ex}
    Cuộc cải cách của Hồ Quý Ly và triều Hồ bắt đầu từ năm nào?
    \shortans{1400}
    \loigiai{Năm 1400, Hồ Quý Ly lên ngôi vua, thiết lập triều Hồ và thực hiện cải cách.}
\end{ex}

%Môn: Toán
\begin{ex}
    Xét tính chất của các hàm số lượng giác:
    \choiceTF
    {\True Hàm số $y = \sin x$ có tập xác định là $\mathbb{R}$.}
    {Hàm số $y = \cos x$ là hàm số lẻ.}
    {\True Hàm số $y = \tan x$ tuần hoàn với chu kì $\pi$.}
    {\True Đồ thị hàm số $y = \sin x$ đi qua gốc tọa độ $O(0;0)$.}
    \loigiai{
        \begin{itemchoice}
        \itemch tập xác định của sin và cos là $\mathbb{R}$.
        \itemch $y = \cos x$ là hàm số chẵn vì $\cos(-x) = \cos x$.
        \itemch chu kì của tan và cot là $\pi$.
        \itemch $\sin(0) = 0$.
        \end{itemchoice}
    }
\end{ex}

%Môn: Vật Lí
\begin{ex}
    Khoảng cách giữa một nút và một bụng sóng liên tiếp trong sóng dừng là $5$ cm. Bước sóng $\lambda$ là bao nhiêu cm?
    \shortans{20}
    \loigiai{Khoảng cách nút - bụng liên tiếp là $\lambda/4 = 5 \Rightarrow \lambda = 20$ cm.}
\end{ex}

%Môn: Sinh Học
\begin{ex}
    Về hệ tuần hoàn ở động vật:
    \choiceTF
    {\True Ở hệ tuần hoàn hở, máu chảy trong động mạch dưới áp lực thấp.}
    {Côn trùng có hệ tuần hoàn kín.}
    {\True Lưỡng cư và bò sát (trừ cá sấu) có tim 3 ngăn.}
    {\True Hệ tuần hoàn kép xuất hiện ở chim và thú.}
    \loigiai{
        \begin{itemchoice}
        \itemch máu trộn lẫn với dịch mô ở xoang cơ thể.
        \itemch côn trùng có hệ tuần hoàn hở.
        \itemch máu bị pha ở tâm thất.
        \itemch gồm vòng tuần hoàn lớn và vòng tuần hoàn nhỏ.
        \end{itemchoice}
    }
\end{ex}

%Môn: Hóa Học
\begin{ex}
    Phân tử Methane ($CH_4$) có bao nhiêu liên kết đơn trong cấu tạo?
    \shortans{4}
    \loigiai{Nguyên tử Carbon trung tâm liên kết với 4 nguyên tử Hydrogen bằng 4 liên kết đơn.}
\end{ex}

%Môn: Giáo Dục Kinh Tế và Pháp Luật
\begin{ex}
    Công dân Việt Nam đủ bao nhiêu tuổi trở lên có quyền tham gia bỏ phiếu bầu cử đại biểu Quốc hội?
    \shortans{18}
    \loigiai{Theo Hiến pháp, mọi công dân đủ 18 tuổi trở lên có quyền bầu cử.}
\end{ex}

%Môn: Toán
\begin{ex}
    Tính đạo hàm của hàm số $y = 3x^2$ tại điểm $x = 2$.
    \shortans{12}
    \loigiai{$y' = 6x$. Tại $x = 2$, $y' = 6 \cdot 2 = 12$.}
\end{ex}

%Môn: Địa Lí
\begin{ex}
    Phát biểu nào sau đây về Liên minh châu Âu (EU) là đúng?
    \choiceTF
    {\True EU là một liên kết kinh tế khu vực lớn nhất thế giới.}
    {Tất cả các nước châu Âu đều là thành viên của EU.}
    {\True Công dân EU có quyền tự do cư trú tại bất cứ nước thành viên nào.}
    {\True Đồng Euro là đồng tiền chung của nhiều quốc gia thành viên EU.}
    \loigiai{
        \begin{itemchoice}
        \itemch dựa trên quy mô GDP và thương mại.
        \itemch ví dụ như Anh đã rời EU, hoặc Thụy Sĩ, Na Uy không phải thành viên.
        \itemch đây là một trong bốn quyền tự do cơ bản.
        \itemch các nước trong khu vực Eurozone.
        \end{itemchoice}
    }
\end{ex}

%Môn: Vật Lí
\begin{ex}
    Cường độ dòng điện $I$ chạy qua điện trở $R = 10 \Omega$ khi hiệu điện thế là $U = 5$ V.
    \shortans{0,5}
    \loigiai{$I = U/R = 5/10 = 0,5$ A.}
\end{ex}

%Môn: Lịch Sử
\begin{ex}
    Ai là người lãnh đạo cuộc khởi nghĩa Lam Sơn (1418 - 1427)?
    \choice
    {Ngô Quyền}
    {Lý Thường Kiệt}
    {\True Lê Lợi}
    {Quang Trung}
    \loigiai{Lê Lợi là người khởi xướng và lãnh đạo nghĩa quân Lam Sơn chống quân Minh.}
\end{ex}

%Môn: Hóa Học
\begin{ex}
    Trong phân tử Nitơ ($N_2$), hai nguyên tử liên kết với nhau bằng bao nhiêu liên kết cộng hóa trị?
    \shortans{3}
    \loigiai{Nitơ có liên kết ba bền vững ($N \equiv N$).}
\end{ex}

%Môn: Địa Lí
\begin{ex}
    \\sochc{2}{
        Đông Nam Á là khu vực năng động, giàu tài nguyên.
    }
    \begin{chc}
        Có bao nhiêu quốc gia nằm trong khu vực này?
        \shortans{11}
        \loigiai{Gồm 10 nước ASEAN và Đông Ti-mo.}
    \end{chc}
    \begin{chc}
        Quốc gia nào ở Đông Nam Á không giáp biển?
        \loigiai{Lào là quốc gia duy nhất không có đường bờ biển.}
    \end{chc}
\end{ex}

%Môn: Sinh Học
\begin{ex}
    Loại hormone thực vật nào kích thích sự nảy mầm của hạt và củ?
    \choice
    {Auxin}
    {\True Gibberellin}
    {Ethylene}
    {Abscisic acid}
    \loigiai{Gibberellin thúc đẩy quá trình phân giải tinh bột thành đường để hạt nảy mầm.}
\end{ex}

%Môn: Toán
\begin{ex}
    Giới hạn $\lim_{x \to 3} (x^2 - 1)$ bằng bao nhiêu?
    \shortans{8}
    \loigiai{$3^2 - 1 = 9 - 1 = 8$.}
\end{ex}

%Môn: Vật Lí
\begin{ex}
    Tụ điện có điện dung $C = 5 \mu F$. Để tích được điện tích $Q = 10 \mu C$ thì cần đặt vào hiệu điện thế $U$ bằng bao nhiêu V?
    \shortans{2}
    \loigiai{$U = Q / C = 10 / 5 = 2$ V.}
\end{ex}

%Môn: Giáo Dục Kinh Tế và Pháp Luật
\begin{ex}
    Khi giá cả của một mặt hàng giảm xuống, lượng cầu của người tiêu dùng đối với mặt hàng đó thường có xu hướng:
    \choice
    {Giảm đi}
    {\True Tăng lên}
    {Không đổi}
    {Biến động không quy luật}
    \loigiai{Theo quy luật cầu, giá giảm thì lượng cầu tăng.}
\end{ex}

%Môn: Hóa Học
\begin{ex}
    Hợp chất $CH_3COOH$ có tên thông thường là gì?
    \choice
    {Rượu etylic}
    {\True Axit axetic}
    {Andehit fomic}
    {Axeton}
    \loigiai{$CH_3COOH$ là thành phần chính của giấm ăn.}
\end{ex}

%Môn: Toán
\begin{ex}
    Trong không gian, cho hai đường thẳng song song $a$ và $b$. Nếu đường thẳng $c$ vuông góc với $a$ thì $c$ có vuông góc với $b$ không?
    \choice
    {Không bao giờ}
    {Chỉ khi $c$ cắt $b$}
    {\True Có}
    {Chỉ khi $c$ chéo $b$}
    \loigiai{Đây là tính chất cơ bản của quan hệ song song và vuông góc.}
\end{ex}

%Môn: Vật Lí
\begin{ex}
    Về hiện tượng cộng hưởng âm:
    \choiceTF
    {\True Xảy ra khi tần số của nguồn âm bằng tần số riêng của hộp đàn.}
    {Làm cho âm thanh nhỏ đi.}
    {\True Giúp tăng cường độ âm.}
    {\True Hộp đàn ghita đóng vai trò là hộp cộng hưởng.}
    \loigiai{
        \begin{itemchoice}
        \itemch 
        \itemch cộng hưởng làm âm to hơn.
        \itemch 
        \itemch 
        \end{itemchoice}
    }
\end{ex}

%Môn: Sinh Học
\begin{ex}
    Số lượng tâm thất trong tim của cá xương là bao nhiêu?
    \shortans{1}
    \loigiai{Tim cá có 2 ngăn: 1 tâm nhĩ và 1 tâm thất.}
\end{ex}

%Môn: Lịch Sử
\begin{ex}
    Chiến thắng nào của quân dân nhà Trần đã kết thúc hoàn toàn mộng xâm lược của quân Nguyên (1288)?
    \choice
    {Đông Bộ Đầu}
    {Chương Dương}
    {Hàm Tử}
    {\True Bạch Đằng}
    \loigiai{Chiến thắng trên sông Bạch Đằng năm 1288 đập tan hoàn toàn ý chí xâm lược của nhà Nguyên.}
\end{ex}

%Môn: Địa Lí
\begin{ex}
    Quốc gia nào sau đây không thuộc khu vực Đông Nam Á?
    \choice
    {Việt Nam}
    {Thái Lan}
    {Lào}
    {\True Ấn Độ}
    \loigiai{Ấn Độ thuộc khu vực Nam Á.}
\end{ex}

%Môn: Hóa Học
\begin{ex}
    Tính khối lượng mol phân tử của nước ($H_2O$). (Lấy $H=1, O=16$).
    \shortans{18}
    \loigiai{$M = 1 \cdot 2 + 16 = 18$.}
\end{ex}

%Môn: Toán
\begin{ex}
    Tính giới hạn $\lim_{n \to \infty} \frac{2n}{n+1}$.
    \shortans{2}
    \loigiai{Chia cả tử và mẫu cho $n$: $\lim \frac{2}{1 + 1/n} = 2$.}
\end{ex}

%Môn: Vật Lí
\begin{ex}
    Một vật dao động tắt dần. Sau mỗi chu kì, năng lượng của nó giảm đi.
    \choiceTF
    {\True Nguyên nhân là do lực ma sát và lực cản môi trường.}
    {Tần số dao động cũng giảm dần theo thời gian.}
    {\True Biên độ dao động giảm dần theo thời gian.}
    {Cơ năng của vật được bảo toàn.}
    \loigiai{
        \begin{itemchoice}
        \itemch ma sát chuyển hóa cơ năng thành nhiệt năng.
        \itemch tần số và chu kì không đổi trong dao động tắt dần (xét mô hình cơ bản).
        \itemch 
        \itemch cơ năng giảm.
        \end{itemchoice}
    }
\end{ex}

%Môn: Lịch Sử
\begin{ex}
    Năm nào thực dân Pháp nổ súng tấn công bán đảo Sơn Trà (Đà Nẵng), mở đầu quá trình xâm lược Việt Nam?
    \shortans{1858}
    \loigiai{Ngày 1-9-1858.}
\end{ex}

%Môn: Sinh Học
\begin{ex}
    Phần lớn thực vật hấp thụ Nitrogen dưới dạng ion nào?
    \choice
    {$N_2$}
    {$NO_2$}
    {\True $NO_3^-$ và $NH_4^+$}
    {$NH_3$}
    \loigiai{Thực vật không hấp thụ được nitơ phân tử trực tiếp mà phải qua dạng muối khoáng.}
\end{ex}

%Môn: Giáo Dục Kinh Tế và Pháp Luật
\begin{ex}
    Về tình trạng thất nghiệp:
    \choiceTF
    {\True Thất nghiệp là tình trạng người lao động muốn làm việc nhưng chưa có việc làm.}
    {Học sinh đang đi học được tính là người thất nghiệp.}
    {\True Thất nghiệp cơ cấu xảy ra khi có sự thay đổi về công nghệ sản xuất.}
    {\True Thất nghiệp gây lãng phí nguồn lực lao động của xã hội.}
    \loigiai{
        \begin{itemchoice}
        \itemch Đúng theo khái niệm kinh tế học.
        \itemch học sinh chưa gia nhập lực lượng lao động.
        \itemch 
        \itemch 
        \end{itemchoice}
    }
\end{ex}

%Môn: Toán
\begin{ex}
    Hàm số $y = \cos x$ đạt giá trị lớn nhất bằng bao nhiêu?
    \shortans{1}
    \loigiai{Giá trị của hàm cos luôn nằm trong đoạn $[-1; 1]$.}
\end{ex}

%Môn: Vật Lí
\begin{ex}
    Sóng điện từ nào có bước sóng ngắn nhất trong các loại sau?
    \choice
    {Sóng vô tuyến}
    {Tia hồng ngoại}
    {Ánh sáng nhìn thấy}
    {\True Tia tử ngoại}
    \loigiai{Theo thang sóng điện từ, tia tử ngoại có tần số cao hơn và bước sóng ngắn hơn 3 loại còn lại.}
\end{ex}

%Môn: Địa Lí
\begin{ex}
    Thủ đô của Nhật Bản là thành phố nào?
    \choice
    {Ky-ô-tô}
    {Ô-xa-ca}
    {\True Tô-ky-ô}
    {Na-gôi-a}
    \loigiai{Tô-ky-ô là trung tâm hành chính, kinh tế của Nhật Bản.}
\end{ex}

%Môn: Hóa Học
\begin{ex}
    Phản ứng nào sau đây là phản ứng đặc trưng của benzene ($C_6H_6$)?
    \choice
    {Chỉ phản ứng cộng}
    {Chỉ phản ứng thế}
    {\True Dễ thế, khó cộng}
    {Dễ cộng, khó thế}
    \loigiai{Do cấu tạo vòng đặc biệt, benzene có tính thơm: dễ thế, khó cộng và bền với tác nhân oxi hóa.}
\end{ex}

%Môn: Toán
\begin{ex}
    Tính đạo hàm của $y = \ln(x)$.
    \shortans{1/x}
    \shortans{0}
    \loigiai{Đây là công thức đạo hàm cơ bản.}
\end{ex}

%Môn: Vật Lí
\begin{ex}
    Đơn vị của mức cường độ âm là gì?
    \choice
    {Oát trên mét vuông ($W/m^2$)}
    {\True Đêxiben ($dB$)}
    {Héc ($Hz$)}
    {Vôn ($V$)}
    \loigiai{Cường độ âm đo bằng $W/m^2$, mức cường độ âm đo bằng $B$ hoặc $dB$.}
\end{ex}

%Môn: Sinh Học
\begin{ex}
    Thành phần nào trong máu đóng vai trò chính trong việc vận chuyển khí Oxy?
    \choice
    {Bạch cầu}
    {Tiểu cầu}
    {\True Hồng cầu (huyết sắc tố)}
    {Huyết tương}
    \loigiai{Hemoglobin trong hồng cầu kết hợp với $O_2$ để vận chuyển đi khắp cơ thể.}
\end{ex}

%Môn: Lịch Sử
\begin{ex}
    Cuộc Cách mạng tháng Mười Nga diễn ra vào năm nào?
    \shortans{1917}
    \loigiai{Cách mạng nổ ra vào ngày 25-10-1917 (theo lịch Nga cũ).}
\end{ex}

%Môn: Toán
\begin{ex}
    Trong một cấp số cộng có $u_1 = 5$ và $d = -2$. Tìm số hạng $u_3$.
    \shortans{1}
    \loigiai{$u_3 = u_1 + 2d = 5 + 2(-2) = 1$.}
\end{ex}

%Môn: Địa Lí
\begin{ex}
    Hoa Kỳ có biên giới phía Nam tiếp giáp với quốc gia nào?
    \choice
    {Canada}
    {Bra-xin}
    {\True Mê-hi-cô}
    {Cu-ba}
    \loigiai{Biên giới phía Nam của Hoa Kỳ giáp với Mê-hi-cô.}
\end{ex}

%Môn: Hóa Học
\begin{ex}
    \\choiceTF
    {\True Ethanol tan vô hạn trong nước.}
    {Phenol có tính acid mạnh hơn acid sulfuric.}
    {\True Ammonia có tính base yếu.}
    {Axit axetic là chất khí ở nhiệt độ thường.}
    \loigiai{
        \begin{itemchoice}
        \itemch Đúng do tạo được liên kết hydrogen với nước.
        \itemch phenol là acid rất yếu, yếu hơn cả acid carbonic.
        \itemch 
        \itemch nó là chất lỏng.
        \end{itemchoice}
    }
\end{ex}

%Môn: Vật Lí
\begin{ex}
    Trong dao động điều hòa, vận tốc biến thiên ______ pha so với li độ.
    \choice
    {Cùng}
    {Ngược}
    {\True Sớm $\pi/2$}
    {Trễ $\pi/2$}
    \loigiai{Vận tốc luôn sớm pha $\pi/2$ so với li độ.}
\end{ex}

%Môn: Sinh Học
\begin{ex}
    Cây xương rồng có lá biến thành gai là một ví dụ về sự thích nghi để giảm bớt quá trình gì?

    \loigiai{Thoát hơi nước.}
\end{ex}

%Môn: Toán
\begin{ex}
    Cho hình chóp $S.ABCD$ có đáy $ABCD$ là hình vuông. Số cặp đường thẳng chéo nhau từ các cạnh của hình chóp là bao nhiêu?
    \loigiai{Sử dụng quy tắc đếm cặp cạnh không cùng nằm trong một mặt phẳng.}
\end{ex}

%Môn: Giáo Dục Kinh Tế và Pháp Luật
\begin{ex}
    Hành vi nào sau đây là vi phạm đạo đức kinh doanh?
    \choice
    {Tăng lương cho nhân viên}
    {Mở rộng thị trường xuất khẩu}
    {\True Xả nước thải độc hại ra sông}
    {Tổ chức hội thảo khách hàng}
    \loigiai{Vi phạm quy định về bảo vệ môi trường là vi phạm đạo đức và pháp luật.}
\end{ex}

%Môn: Hóa Học
\begin{ex}
    Sản phẩm chính của phản ứng giữa Benzene và Chlorine (xúc tác bột sắt) là gì?

    \loigiai{Chlorobenzene ($C_6H_5Cl$).}
\end{ex}

%Môn: Địa Lí
\begin{ex}
    Về dân cư khu vực Mỹ La-tinh:
    \choiceTF
    {\True Là khu vực có sự đa dạng về chủng tộc bậc nhất thế giới.}
    {Tỉ lệ dân thành thị của Mỹ La-tinh rất thấp.}
    {\True Có sự chênh lệch giàu nghèo rất lớn trong xã hội.}
    {\True Dân cư tập trung đông ở vùng ven biển.}
    \loigiai{
        \begin{itemchoice}
        \itemch do quá trình nhập cư lâu đời.
        \itemch tỉ lệ dân thành thị ở đây rất cao (khoảng 80%).
        \itemch 
        \itemch 
        \end{itemchoice}
    }
\end{ex}

%Môn: Lịch Sử
\begin{ex}
    Dưới thời vua Minh Mạng, cơ quan nào làm nhiệm vụ tư vấn cao nhất cho nhà vua về việc quân sự?
    \choice
    {Nội các}
    {\True Cơ mật viện}
    {Quốc tử giám}
    {Hàn lâm viện}
    \loigiai{Cơ mật viện được thành lập năm 1834 để giải quyết các việc cơ mật.}
\end{ex}

%Môn: Toán
\begin{ex}
    Tính tổng của cấp số nhân lùi vô hạn $1; 1/3; 1/9; ...$
    \shortans{1,5}
    \loigiai{$S = u_1 / (1-q) = 1 / (1 - 1/3) = 1 / (2/3) = 1,5$.}
\end{ex}

%Môn: Vật Lí
\begin{ex}
    Trong thí nghiệm Young về giao thoa ánh sáng, khoảng cách giữa hai khe là $a = 2$ mm, khoảng cách đến màn $D = 2$ m, bước sóng $\lambda = 0,5 \mu m$. Tính khoảng vân $i$ (mm).
    \shortans{0,5}
    \loigiai{$i = \lambda D / a = (0,5 \cdot 10^{-3} \cdot 2000) / 2 = 0,5$ mm.}
\end{ex}

%Môn: Sinh Học
\begin{ex}
    Phát biểu về quá trình tiêu hóa:
    \choiceTF
    {\True Tiêu hóa ở dạ dày người bao gồm cả biến đổi cơ học và hóa học.}
    {Ruột non là nơi diễn ra sự hấp thụ nước chính yếu.}
    {\True Enzyme Amylase trong nước bọt giúp tiêu hóa tinh bột chín.}
    {Mật do túi mật sản xuất ra.}
    \loigiai{
        \begin{itemchoice}
        \itemch 
        \itemch ruột già mới là nơi hấp thụ nước chủ yếu, ruột non hấp thụ dinh dưỡng.
        \itemch 
        \itemch gan sản xuất mật, túi mật chỉ dự trữ.
        \end{itemchoice}
    }
\end{ex}

%Môn: Toán
\begin{ex}
    Tìm nghiệm của phương trình $\log_5(x) = 2$.
    \shortans{25}
    \loigiai{$x = 5^2 = 25$.}
\end{ex}

%Môn: Vật Lí
\begin{ex}
    \\immini[thm]{
        Một sóng ngang truyền trên một sợi dây. Tại một thời điểm, hình dạng sợi dây như hình vẽ. Hướng truyền sóng là từ trái sang phải. Điểm M trên dây đang chuyển động theo hướng nào?
    }{
        [Image of a wave moving right, M is on the rising slope]}
    \choice
    {\True Đi lên}
    {Đi xuống}
    {Sang trái}
    {Sang phải}
    \loigiai{Dựa vào phương pháp "đỉnh sóng đuổi nhau", điểm M nằm trước đỉnh sóng đang tới nên nó sẽ đi lên.}
\end{ex}

%Môn: Địa Lí
\begin{ex}
    Năm 2020, quốc gia nào sau đây có chỉ số HDI cao nhất khu vực Đông Nam Á?

    \loigiai{Xin-ga-po (Singapore).}
\end{ex}

%Môn: Hóa Học
\begin{ex}
    Andehit fomic ($HCHO$) có thể tham gia phản ứng tráng bạc. Một phân tử $HCHO$ tạo ra tối đa bao nhiêu nguyên tử Bạc ($Ag$)?
    \shortans{4}
    \loigiai{Riêng $HCHO$ khi phản ứng hoàn toàn tạo ra $4Ag$.}
\end{ex}

%Môn: Giáo Dục Kinh Tế và Pháp Luật
\begin{ex}
    Quyền bất khả xâm phạm về thân thể có nghĩa là không ai bị bắt nếu không có quyết định của Tòa án hoặc phê chuẩn của ______.

    \loigiai{Viện kiểm sát.}
\end{ex}

%Môn: Lịch Sử
\begin{ex}
    Cách mạng tư sản Anh (thế kỉ XVII) đã thiết lập chế độ chính trị nào?
    \choice
    {Cộng hòa}
    {Quân chủ chuyên chế}
    {\True Quân chủ lập hiến}
    {Xã hội chủ nghĩa}
    \loigiai{Sau thời kỳ biến động, Anh duy trì ngôi vua nhưng quyền lực thuộc về Quốc hội.}
\end{ex}

%Môn: Toán
\begin{ex}
    Tính giới hạn $\lim_{x \to 0} \frac{x^2 - x}{x}$.
    \shortans{-1}
    \loigiai{$\lim_{x \to 0} (x - 1) = -1$.}
\end{ex}

%Môn: Vật Lí
\begin{ex}
    Đơn vị đo tần số dao động là gì?
    \shortans{Hz}
    \shortans{0}
    \loigiai{Héc (Hz).}
\end{ex}

%Môn: Sinh Học
\begin{ex}
    Nhóm thực vật nào sau đây thực hiện chu trình CAM?
    \choice
    {Lúa, ngô}
    {Xoài, nhãn}
    {\True Dứa, thanh long, xương rồng}
    {Rêu, dương xỉ}
    \loigiai{Thực vật CAM sống ở nơi khô hạn, đóng khí khổng vào ban ngày.}
\end{ex}

%Môn: Hóa Học
\begin{ex}
    Khi cho mẩu Natri (Sodium) vào Ethanol, có khí gì thoát ra?
    \shortans{H2}
    \loigiai{Phản ứng giải phóng khí Hydrogen ($H_2$).}
\end{ex}

%Môn: Toán
\begin{ex}
    \\sochc{2}{
        Cho hình lập phương $ABCD.A'B'C'D'$ cạnh $1$.
    }
    \begin{chc}
        Độ dài đường chéo $AC$ của mặt đáy là bao nhiêu? (Lấy $\sqrt{2} \approx 1,41$).
        \shortans{1,41}
        \loigiai{$AC = \sqrt{1^2 + 1^2} = \sqrt{2} \approx 1,41$.}
    \end{chc}
    \begin{chc}
        Đường thẳng $AA'$ có vuông góc với mặt phẳng $(ABCD)$ không?
        \loigiai{Có, vì đây là hình lập phương.}
    \end{chc}
\end{ex}

%Môn: Địa Lí
\begin{ex}
    Phần lớn lãnh thổ Liên bang Nga nằm trong đới khí hậu nào?
    \choice
    {Nhiệt đới}
    {\True Ôn đới}
    {Cận nhiệt đới}
    {Xích đạo}
    \loigiai{Nga có khí hậu ôn đới lục địa chiếm diện tích lớn nhất.}
\end{ex}

%Môn: Sinh Học
\begin{ex}
    Về sự sinh trưởng ở động vật:
    \choiceTF
    {\True Phát triển qua biến thái hoàn toàn có giai đoạn nhộng.}
    {Châu chấu phát triển qua biến thái hoàn toàn.}
    {\True Ở người, phát triển không qua biến thái.}
    {\True Hormone sinh trưởng (GH) do tuyến yên tiết ra.}
    \loigiai{
        \begin{itemchoice}
        \itemch ví dụ như bướm, ruồi.
        \itemch châu chấu biến thái không hoàn toàn.
        \itemch con non sinh ra có hình dạng giống người lớn.
        \itemch 
        \end{itemchoice}
    }
\end{ex}

%Môn: Vật Lí
\begin{ex}
    Một vật dao động điều hòa với biên độ $A = 10$ cm. Quãng đường vật đi được trong một chu kì là bao nhiêu cm?
    \shortans{40}
    \loigiai{$S = 4A = 4 \cdot 10 = 40$ cm.}
\end{ex}

%Môn: Hóa Học
\begin{ex}
    Phân tử Axetilen ($C_2H_2$) có bao nhiêu liên kết pi ($\pi$)?
    \shortans{2}
    \loigiai{Liên kết ba $C \equiv C$ gồm 1 liên kết sigma và 2 liên kết pi.}
\end{ex}

%Môn: Toán
\begin{ex}
    Cho hàm số $y = x^2 - 4x + 3$. Đạo hàm $y' = 0$ tại $x$ bằng bao nhiêu?
    \shortans{2}
    \loigiai{$y' = 2x - 4$. $y' = 0 \Leftrightarrow x = 2$.}
\end{ex}

%Môn: Lịch Sử
\begin{ex}
    Cuộc chiến tranh bảo vệ Tổ quốc đầu tiên của nhân dân ta chống lại quân xâm lược phương Bắc là cuộc kháng chiến chống quân nào?
    \choice
    {Tống}
    {\True Tần}
    {Hán}
    {Nguyên}
    \loigiai{Cuộc kháng chiến chống quân Tần của người Âu Việt và Lạc Việt.}
\end{ex}

%Môn: Giáo Dục Kinh Tế và Pháp Luật
\begin{ex}
    Quyền được đảm bảo an toàn và bí mật thư tín, điện thoại là quyền của ai?
    \choice
    {Chỉ công chức}
    {Chỉ người có quốc tịch}
    {\True Mọi người}
    {Chỉ người lớn}
    \loigiai{Đây là quyền cơ bản của con người được ghi trong Hiến pháp.}
\end{ex}

%Môn: Địa Lí
\begin{ex}
    Quốc gia nào sau đây có dân số già nhất thế giới?
    \choice
    {Việt Nam}
    {Hoa Kỳ}
    {\True Nhật Bản}
    {Bra-xin}
    \loigiai{Nhật Bản đang đối mặt với tình trạng già hóa dân số nghiêm trọng nhất.}
\end{ex}

%Môn: Toán
\begin{ex}
    Tính $\log_{10}(1000)$.
    \shortans{3}
    \loigiai{$10^3 = 1000$.}
\end{ex}

%Môn: Vật Lí
\begin{ex}
    Trong mạch điện, linh kiện nào dùng để ngăn dòng điện một chiều và cho dòng điện xoay chiều đi qua?
    \choice
    {Điện trở}
    {Cuộn dây}
    {\True Tụ điện}
    {Ampe kế}
    \loigiai{Đây là tính chất cơ bản của tụ điện.}
\end{ex}

%Môn: Hóa Học
\begin{ex}
    Sản phẩm của phản ứng thủy phân Ester trong môi trường kiềm được gọi là phản ứng gì?

    \loigiai{Xà phòng hóa.}
\end{ex}

%Môn: Sinh Học
\begin{ex}
    Cấu tạo của một Neuron gồm thân, sợi nhánh và ______.

    \loigiai{Sợi trục (Axon).}
\end{ex}

%Môn: Toán
\begin{ex}
    Số hạng tổng quát của dãy số lẻ $1, 3, 5, 7, ...$ là $u_n = 2n - k$. Tìm $k$.
    \shortans{1}
    \loigiai{$u_1 = 2(1) - k = 1 \Rightarrow k = 1$.}
\end{ex}

%Môn: Địa Lí
\begin{ex}
    Tây Nam Á giàu có nhờ tài nguyên nào là chủ yếu?
    \choice
    {Vàng}
    {Kim cương}
    {\True Dầu mỏ}
    {Than đá}
    \loigiai{Dầu mỏ là nguồn thu nhập chính của khu vực này.}
\end{ex}

%Môn: Vật Lí
\begin{ex}
    Khi sóng truyền từ không khí vào nước, bước sóng $\lambda$ sẽ như thế nào?
    \choice
    {\True Tăng lên}
    {Giảm đi}
    {Không đổi}
    {Bằng 0}
    \loigiai{Vì $v$ tăng ($v_{nuoc} > v_{khi}$) và $f$ không đổi nên $\lambda = v/f$ tăng.}
\end{ex}

%Môn: Hóa Học
\begin{ex}
    Hợp chất nào sau đây là hydrocarbon thơm?
    \choice
    {Methane}
    {Ethylene}
    {Acetylene}
    {\True Benzene}
    \loigiai{Benzene là đại diện tiêu biểu của hydrocarbon thơm.}
\end{ex}

%Môn: Toán
\begin{ex}
    Tính đạo hàm của $y = \tan(x)$ tại $x=0$.
    \shortans{1}
    \loigiai{$y' = 1/\cos^2(x)$. Tại $x=0$, $y' = 1/1 = 1$.}
\end{ex}

%Môn: Lịch Sử
\begin{ex}
    Nhà nước Liên bang Cộng hòa xã hội chủ nghĩa Xô viết viết tắt là gì?
    \loigiai{Liên Xô (LX).}
\end{ex}

%Môn: Địa Lí
\begin{ex}
    Chỉ số GNI/người của Việt Nam năm 2020 là bao nhiêu USD? (Làm tròn số nguyên).
    \shortans{3390}
    \loigiai{Theo bảng 1.1 trang 8.}
\end{ex}

%Môn: Sinh Học
\begin{ex}
    \\choiceTF
    {\True Thực vật C3 là nhóm thực vật phổ biến nhất.}
    {Cây dứa là thực vật C4.}
    {\True Quá trình hô hấp ở thực vật giải phóng năng lượng dưới dạng ATP.}
    {\True Nước là nguyên liệu của quá trình quang hợp.}
    \loigiai{
        \begin{itemchoice}
        \itemch 
        \itemch dứa là thực vật CAM.
        \itemch 
        \itemch 
        \end{itemchoice}
    }
\end{ex}

%Môn: Vật Lí
\begin{ex}
    Một con lắc lò xo có độ cứng $k = 100 N/m$, vật nặng $m = 1 kg$. Tần số góc $\omega$ là bao nhiêu rad/s?
    \shortans{10}
    \loigiai{$\omega = \sqrt{k/m} = \sqrt{100/1} = 10$ rad/s.}
\end{ex}

%Môn: Toán
\begin{ex}
    Tính giới hạn $\lim_{x \to \infty} (1/x + 2)$.
    \shortans{2}
    \loigiai{$0 + 2 = 2$.}
\end{ex}

%Môn: Địa Lí
\begin{ex}
    Tỉ lệ dân thành thị của Hoa Kỳ năm 2020 là bao nhiêu \%? (Lấy số nguyên).
    \shortans{82}
    \loigiai{Theo SGK là 82,7\%, làm tròn là 83 hoặc lấy theo bảng là 82.}
\end{ex}

%Môn: Hóa Học
\begin{ex}
    Dung dịch chất nào sau đây có pH < 7?
    \choice
    {$NaOH$}
    {$NaCl$}
    {\True $HCl$}
    {$NH_3$}
    \loigiai{Acid có pH < 7.}
\end{ex}

%Môn: Sinh Học
\begin{ex}
    Hiện tượng thụ phấn ở thực vật là sự vận chuyển hạt phấn đến bộ phận nào của hoa?

    \loigiai{Núm nhụy.}
\end{ex}

%Môn: Lịch Sử
\begin{ex}
    Năm nào cuộc khởi nghĩa của Triệu Thị Trinh bùng nổ?
    \shortans{248}
    \loigiai{Năm 248, bà Triệu khởi nghĩa chống quân Ngô.}
\end{ex}

%Môn: Giáo Dục Kinh Tế và Pháp Luật
\begin{ex}
    Hành vi tự ý bóc thư của người khác là vi phạm quyền ______ của công dân.

    \loigiai{Quyền được bảo đảm an toàn và bí mật thư tín.}
\end{ex}

%Môn: Toán
\begin{ex}
    Cho hình chóp $S.ABC$. Có bao nhiêu cặp đường thẳng chéo nhau tạo bởi các cạnh bên và cạnh đáy?
    \shortans{3}
    \loigiai{Các cặp: (SA, BC), (SB, AC), (SC, AB).}
\end{ex}